\section*{अध्याय \ref{ch:g12:age-of-the-earth} --- रेडियोमितीय काल-निर्धारण और पृथ्वी की आयु}

\begin{solution}{exo:g12:age-of-the-earth:1}
(a) \omterm{def:g12:age-of-the-earth:closed-system}{बंद तंत्र}; (b) \omterm{def:g12:age-of-the-earth:isochron}{समकालरेखा}; (c) \omterm{def:g12:age-of-the-earth:age}{पृथ्वी की आयु}।
\end{solution}

\begin{solution}{exo:g12:age-of-the-earth:2}
प्लेजियोक्लेज़ (0.35), हॉर्नब्लेंड (1.20), पूर्ण चट्टान (3.10), K-फ़ेल्डस्पार (6.40),
बायोटाइट (24.0)। बायोटाइट सबसे ऊँचा पड़ता है: उसमें अपने स्ट्रॉन्शियम-86 की तुलना में सबसे अधिक
रुबिडियम-87 है, इसलिए ग्रेनाइट के बनने के बाद से उसके रुबिडियम ने सबसे अधिक
स्ट्रॉन्शियम-87 बनाया है।
\end{solution}

\begin{solution}{exo:g12:age-of-the-earth:3}
$T = \ln 2/\lambda = 0.6931/\num{1.3972e-11} \approx \num{4.96e10}$ वर्ष,
यानी 49.6 अरब वर्ष: \omterm{def:g12:age-of-the-earth:age}{पृथ्वी की आयु} का $49.6/4.54 \approx 11$ गुना।
\end{solution}

\begin{solution}{exo:g12:age-of-the-earth:4}
$\lambda t = \num{1.3972e-11} \times \num{3.10e8} \approx 0.004331$, और
$e^{0.004331} - 1 \approx 0.004341$। 480 मिलियन वर्ष के लिए, $\lambda t
\approx 0.006707$ और $e^{\lambda t} - 1 \approx 0.006729$। दोनों स्वयं
$\lambda t$ के निकट हैं, क्योंकि $\lambda t$ छोटा है।
\end{solution}

\begin{solution}{exo:g12:age-of-the-earth:5}
फ़िलिप्स (1860), केल्विन (1862), जॉली (1899), होम्स (1927), पैटरसन
(1956)। $4540/100 \approx 45$ गुना।
\end{solution}

\begin{solution}{exo:g12:age-of-the-earth:6}
आज $0.7080 + 10.0 \times 0.004341 \approx 0.7514$; ग्रेनाइट के बनने के समय 0.7080,
ग्रेनाइट के हर खनिज की तरह।
\end{solution}

\begin{solution}{exo:g12:age-of-the-earth:7}
(1) $m = (0.77525 - 0.71806)/(9.40 - 0.90) = 0.05719/8.50 \approx 0.006728$;
$t = \ln(1.006728)/\num{1.3972e-11} \approx \num{4.80e8}$ वर्ष, यानी 480 मिलियन
वर्ष। (2) $0.71806 + 1.70 \times 0.006728 \approx 0.72950$ और $0.71806 + 4.20
\times 0.006728 \approx 0.74632$: दोनों रेखा पर हैं। (3) हाँ: ग्रेनाइट
शिस्ट को काटता है, इसलिए वह नया है, और 310, 480 मिलियन वर्ष से कम है।
\end{solution}

\begin{solution}{exo:g12:age-of-the-earth:8}
उसके सभी बिंदु एक ही स्थान पर पड़ेंगे: एक बिंदु न कोई रेखा तय करता है, न
कोई ढाल। बायोटाइट दूसरों से बहुत दाईं ओर पड़ता है; तब किसी भी बिंदु की एक छोटी
त्रुटि उससे और प्लेजियोक्लेज़ से होकर जाने वाली रेखा के ढाल को उससे कहीं कम बदलती है
जितना वह पास-पास के दो बिंदुओं के लिए बदलती।
\end{solution}

\begin{solution}{exo:g12:age-of-the-earth:9}
(1) उसका रुबिडियम नहीं बदला है, इसलिए उसका $x$ वही रहता है; उसमें
स्ट्रॉन्शियम-87 कम है, इसलिए उसका $y$ गिरता है: बिंदु सीधे नीचे खिसकता है। (2) उससे और
प्लेजियोक्लेज़ से होकर जाने वाली रेखा कम ढलवाँ होती है, और बहुत कम आयु देती है।
(3) बायोटाइट अब दूसरे खनिजों से होकर जाने वाली रेखा पर नहीं पड़ता: बिंदु
एक रेखा में नहीं हैं, जो इसका संकेत है कि एक नमूना \omterm{def:g12:age-of-the-earth:closed-system}{बंद तंत्र} नहीं रहा है।
\end{solution}

\begin{solution}{exo:g12:age-of-the-earth:10}
(1) $4540/96 \approx 47$ गुना। (2) यदि कुछ नमक समुद्र से निकल जाता है, उदाहरण के लिए
अवसादों में दबकर, तो समुद्र में उस सारे नमक से कम नमक होता है जो नदियाँ लाईं:
मौजूद नमक को वार्षिक आपूर्ति से भाग देने पर बहुत कम
समय मिलता है।
\end{solution}

\begin{solution}{exo:g12:age-of-the-earth:11}
$m = (34.86 - 10.34)/(50.28 - 9.46) = 24.52/40.82 \approx 0.601$, जो
4.55 (0.590) और 4.60 अरब वर्ष (0.610) के ढालों के बीच है: लगभग 4.58
अरब वर्ष।
\end{solution}

\begin{solution}{exo:g12:age-of-the-earth:12}
(1) बेसाल्ट, कम अनुपात वाला, प्रावार से; ग्रेनाइट, ऊँचे अनुपात
वाला, कम से कम आंशिक रूप से पिघली भूपर्पटी से। (2) प्रावार के पिघलने पर रुबिडियम
गलित में चला जाता है, इसलिए वह भूपर्पटी में सांद्रित होता गया है, जिसने
युगों में अधिक स्ट्रॉन्शियम-87 बनाया है। (3) नहीं: एक ही मैग्मा से क्रिस्टलित होने वाले
खनिज उसका प्रारंभिक अनुपात साझा करते हैं, और ये दोनों चट्टानें नहीं करतीं।
\end{solution}

\begin{solution}{exo:g12:age-of-the-earth:13}
(1) रुबिडियम-87: $2^{-4.54/49.61} \approx 0.94$, लगभग \qty{94}{\percent}
बचा है; यूरेनियम-238: $2^{-4.54/4.47} \approx 0.49$, लगभग आधा। (2) $m =
e^{\lambda t} - 1 \approx \num{1.3972e-11} \times 10^6 \approx
\num{1.4e-5}$: वार्डन ग्रेनाइट का बायोटाइट भी ($x = 24$) प्रारंभिक अनुपात से
केवल $0.0003$ ऊपर उठेगा। इतना कम स्ट्रॉन्शियम-87 बना है कि ढाल
को शून्य से अलग पहचानना कठिन है: धीमी घड़ी पुरानी चट्टानों के लिए उपयुक्त है।
\end{solution}

\begin{solution}{exo:g12:age-of-the-earth:14}
(1) $4.438 + 0.6026 \times 19.0 \approx 15.89$; $15.5$ इससे $0.39$ कम है,
लगभग \qty{2.4}{\percent}: अवसाद रेखा के निकट पड़ता है। (2) 0.6026,
0.590 (4.55) और 0.610 (4.60) के बीच है: लगभग 4.58 अरब वर्ष, पैटरसन के
4.48 से 4.62 अरब वर्ष के परास के भीतर।
\end{solution}

\begin{solution}{exo:g12:age-of-the-earth:15}
(1) तब भंडार शुद्ध रूप से केवल $46.7 - 24.4 = \qty{22.3}{TW}$ खोता है, इसलिए वही
भंडार $46.7/22.3 \approx 2.1$ गुना अधिक समय तक चलता है। (2) $96 \times 2.1 \approx
200$ मिलियन वर्ष, अब भी बीस गुने से अधिक कम। (3) \omterm{def:g12:internal-heat:convection}{प्रावार
संवहन}, जो ऊष्मा को गहरे भीतर से ऊपर लाता है और सतह के निकट प्रवणता को
केवल चालन की अपेक्षा कहीं अधिक समय तक ऊँचा रखता है।
\end{solution}

\begin{solution}{pb:g12:age-of-the-earth:1}
\textbf{1.} बायोटाइट ($x = 24.0$)।

\textbf{2.} वे एक ही मैग्मा से क्रिस्टलित हुए, और एक ही गलित से क्रिस्टलित होने वाले
खनिज उसका स्ट्रॉन्शियम उसके समस्थानिक-अनुपात के साथ ग्रहण करते हैं।

\textbf{3.} स्ट्रॉन्शियम-86 स्थायी है और किसी क्षय से नहीं बनता: \omterm{def:g12:age-of-the-earth:closed-system}{बंद तंत्र} में
उसकी संख्या नहीं बदलती, इसलिए वह एक स्थिर मापदंड है।

\textbf{4.} स्ट्रॉन्शियम-87, जो रुबिडियम-87 के क्षय से बनता है; रुबिडियम-87
घटा है और स्ट्रॉन्शियम-86 वैसा ही रहा है।

\textbf{5.} $y = 0.7080$ पर एक क्षैतिज रेखा: अनुपात वही, $x$ भिन्न।

\textbf{6.} $(0.81218 - 0.70952)/(24.0 - 0.35) = 0.10266/23.65 \approx
0.004341$।

\textbf{7.} $(0.72146 - 0.70952)/(3.10 - 0.35) = 0.01194/2.75 \approx
0.004342$।

\textbf{8.} वे 4000 में 1 के भीतर मेल खाते हैं। पहला अधिक विश्वसनीय है:
उसके दोनों बिंदु दूर-दूर हैं, इसलिए अनुपातों का पूर्णांकन कम मायने रखता है।

\textbf{9.} $0.70952 - 0.35 \times 0.004341 \approx 0.7080$।

\textbf{10.} $0.7080 + 6.40 \times 0.004341 \approx 0.73578$: ठीक
मापा गया मान।

\textbf{11.} $e^{\lambda t} = 1 + m = 1.004341$।

\textbf{12.} $\lambda t = \ln 1.004341 \approx 0.004331$।

\textbf{13.} $t = 0.004331/\num{1.3972e-11} \approx \num{3.10e8}$ वर्ष।

\textbf{14.} हाँ: ग्रेनाइट शिस्ट को काटता है, इसलिए अनुप्रस्थ काट के
सिद्धांत से वह नया है, और 310, 480 मिलियन
वर्ष से कम है।

\textbf{15.} 310 मिलियन वर्ष, 358.9 और 298.9 मिलियन वर्ष के बीच पड़ता है:
कार्बोनिफ़ेरस।

\textbf{16.} ग्रेनाइट के क्रिस्टलित होने के समय मैग्मा का अनुपात
$^{87}\mathrm{Sr}/^{86}\mathrm{Sr}$।

\textbf{17.} ग्रेनाइट: $0.7080 > 0.7035$।

\textbf{18.} प्रावार के पिघलने पर रुबिडियम गलित में चला जाता है, इसलिए वह
भूपर्पटी में सांद्रित होता गया है; अधिक रुबिडियम के साथ, भूपर्पटी ने अधिक
स्ट्रॉन्शियम-87 बनाया है।

\textbf{19.} यह कि मैग्मा, कम से कम आंशिक रूप से, केवल प्रावार से नहीं, बल्कि पिघली
\omterm{def:g12:continental-crust:continental-crust}{महाद्वीपीय भूपर्पटी} से आया।

\textbf{20.} वार्डन ग्रेनाइट \textbf{310 मिलियन वर्ष} पुराना है।
\end{solution}
